% ═════════════════════════════════════════════════════════════
% Problem Properties
\label{sec:properties}
% ═════════════════════════════════════════════════════════════

\noindent Before discussing the optimization algorithms, we analyze the mathematical
properties of the objective function~\eqref{eq:expanded} that are relevant for
convergence.


% ─────────────────────────────────────────────────────────────
\subsection{Gradient and Hessian}
% ─────────────────────────────────────────────────────────────

\noindent The objective~\eqref{eq:expanded} is a quadratic function of $w$.
Consequently, its Hessian matrix is constant, as it does not depend on the
variable $w$ but only on the fixed data matrix $X$ and the regularization
parameter $\lambda$. Its gradient and Hessian are:

\begin{align}
    \nabla f(w) &= X(X^T w - y) + \lambda^2 w,
    \label{eq:gradient} \\
    H \coloneqq \nabla^2 f(w) &= XX^T + \lambda^2 I_m.
    \label{eq:hessian}
\end{align}


% ─────────────────────────────────────────────────────────────
\subsection{Strong Convexity and Smoothness}
% ─────────────────────────────────────────────────────────────

\noindent To guarantee the existence of a unique global minimum and to analyze the
convergence speed of our optimization algorithms, we first establish the strong
convexity and smoothness of the objective function.

\noindent By Lemma~\ref{lemma:fullrank}, the Hessian $H = XX^T + \lambda^2 I_m$
is symmetric positive definite whenever $\lambda > 0$. This confirms that $f(w)$
is strictly convex, ensuring that any local minimum is also the unique global
minimum.

\noindent Formally, this implies that $f$ is $\mu$-strongly convex and $L$-smooth,
where the constants $\mu$ and $L$ correspond respectively to the smallest and
largest eigenvalues of $H$. By Lemma~\ref{lemma:eigenvalues} (eigenvalue shift),
the eigenvalues of $H = XX^T + \lambda^2 I_m$ are obtained by shifting those of
$XX^T$ by $\lambda^2$:
\begin{equation}
\label{eq:muL}
    \mu = \lambda_{\min}(H) = \lambda_{\min}(XX^T) + \lambda^2,
    \qquad
    L   = \lambda_{\max}(H) = \lambda_{\max}(XX^T) + \lambda^2.
\end{equation}

\noindent The smoothness constant $L$ bounds how rapidly the gradient can change.
For all $w, w' \in \mathbb{R}^m$:
\begin{equation}
    \norm{\nabla f(w) - \nabla f(w')} = \norm{H(w - w')} \leq \norm{H}\,\norm{w - w'} = L\norm{w - w'}.
\end{equation}

\noindent Since our data matrix $X \in \mathbb{R}^{m \times n}$ is tall-thin
($m > n$), by Lemma~\ref{lemma:singvals} the eigenvalues of $XX^T$ are
$\{\sigma_1^2, \ldots, \sigma_n^2, 0, \ldots, 0\}$, so in particular
$\lambda_{\min}(XX^T) = 0$. Therefore, the strong convexity constant simplifies
to:
\begin{equation}
    \mu = \lambda^2.
\end{equation}
This highlights the crucial role of $\lambda$: without regularization, $H$ would
be only positive \emph{semi}-definite, $\mu = 0$, and $f$ would lack strong
convexity. In that case, the system $Hw = Xy$ would be underdetermined, the
minimum would not be unique, and iterative methods would not converge to a
well-defined solution.

\noindent Together, $L$ and $\mu$ quantify the curvature of the objective: $\mu$
controls the minimum curvature (how sharply $f$ rises in the flattest direction)
and $L$ the maximum curvature (how sharply it rises in the steepest direction).
Their ratio defines the condition number of the problem, which directly impacts
the convergence speed of iterative solvers such as L-BFGS.


% ─────────────────────────────────────────────────────────────
\subsection{Condition Number}
\label{sec:conditioning}
% ─────────────────────────────────────────────────────────────

\noindent The geometric shape of the objective function---and consequently the
difficulty of the optimization problem---is summarized by the condition number of
the augmented matrix $\hat{X}$, defined as the ratio of its largest to smallest
singular values.

\noindent To express this in terms of the problem parameters, we note that the
singular values of $\hat{X}$ are the square roots of the eigenvalues of
$\hat{X}^T\hat{X}$. A direct computation gives:
\begin{equation}
\label{eq:Xhat_gram}
    \hat{X}^T\hat{X}
    = \begin{bmatrix} X & \lambda I_m \end{bmatrix}
      \begin{bmatrix} X^T \\ \lambda I_m \end{bmatrix}
    = XX^T + \lambda^2 I_m = H,
\end{equation}
so $\sigma_i(\hat{X}) = \sqrt{\lambda_i(H)}$ for each $i$. Combining this with
equations~\eqref{eq:muL}, the condition number of $\hat{X}$ is:
\begin{equation}
\label{eq:condition}
    \kappa(\hat{X})
    = \frac{\sigma_{\max}(\hat{X})}{\sigma_{\min}(\hat{X})}
    = \sqrt{\frac{\lambda_{\max}(H)}{\lambda_{\min}(H)}}
    = \sqrt{\frac{L}{\mu}}
    = \frac{\sqrt{\lambda_{\max}(XX^T) + \lambda^2}}{\lambda}.
\end{equation}

\noindent As the equation shows, the condition number depends heavily on the
regularization parameter. For small values of $\lambda$, it scales as
$\mathcal{O}(1/\lambda)$: weak regularization produces a highly elongated,
canyon-like objective landscape, causing iterative methods to take many small
zigzagging steps. Conversely, a large $\lambda$ drives $\kappa(\hat{X}) \to 1$,
making the problem nearly spherical and easy to solve from any direction.